% =============================================================================
% hpc_complete_learning_guide.tex
% Master document — includes all six topic files
% EC7207 High Performance Computing – Group 11
%
% Compile with:
%   pdflatex hpc_complete_learning_guide.tex   (run twice for TOC)
% =============================================================================
\documentclass[12pt,a4paper]{report}

\usepackage[margin=2.5cm]{geometry}
\usepackage{listings}
\usepackage{xcolor}
\usepackage{amsmath}
\usepackage{amssymb}
\usepackage{booktabs}
\usepackage[colorlinks=true,linkcolor=blue,urlcolor=blue]{hyperref}
\usepackage{tikz}
\usepackage{array}
\usepackage{enumitem}
\usepackage{fancyhdr}
\usepackage{titlesec}
\usepackage{graphicx}
\usepackage{tcolorbox}
\usepackage{multicol}

% ── Code listing style ────────────────────────────────────────────────────────
\lstdefinestyle{cstyle}{
    language=C,
    basicstyle=\ttfamily\small,
    keywordstyle=\color{blue}\bfseries,
    commentstyle=\color{gray}\itshape,
    stringstyle=\color{red!70!black},
    numbers=left, numberstyle=\tiny\color{gray},
    numbersep=8pt, frame=single, rulecolor=\color{black!30},
    breaklines=true, showstringspaces=false,
    backgroundcolor=\color{gray!5},
    tabsize=4
}
\lstset{style=cstyle}

\pagestyle{fancy}
\fancyhf{}
\lhead{EC7207 HPC – Group 11}
\rhead{Complete Learning Guide}
\cfoot{\thepage}

% =============================================================================
\begin{document}
% =============================================================================

% ── Title Page ────────────────────────────────────────────────────────────────
\begin{titlepage}
\centering
\vspace*{3cm}
{\Huge\bfseries HPC Learning Guide\\[0.5em]
Pearson Correlation\\[0.2em]
Recommender System}\\[2em]
{\Large EC7207 High Performance Computing}\\[1em]
{\large Group 11}\\[3em]
\begin{tcolorbox}[colback=blue!5,colframe=blue!50,width=0.85\textwidth]
\centering
\textbf{Parallelisation Versions Covered:}\\[0.8em]
\begin{tabular}{lll}
1. & Serial Baseline & Algorithm foundation \\
2. & OpenMP & Shared-memory, fork-join \\
3. & Pthreads & Manual thread management \\
4. & MPI & Distributed-memory, message passing \\
5. & CUDA & GPU massively parallel \\
6. & Hybrid MPI+OpenMP & Two-level parallelism \\
\end{tabular}
\end{tcolorbox}
\vfill
{\large\today}
\end{titlepage}

% ── Table of Contents ─────────────────────────────────────────────────────────
\tableofcontents
\newpage

% =============================================================================
\chapter*{Introduction}
\addcontentsline{toc}{chapter}{Introduction}
% =============================================================================

\section*{Project Overview}

This project implements a \textbf{User-Based Collaborative Filtering
Recommender System} using Pearson Correlation Similarity, parallelised with
five different HPC techniques.

\textbf{The algorithm in one sentence:}  Given a sparse matrix of user ratings
($1$--$5$ stars), predict what rating each user would give to items they
haven't seen, using the ratings of the most similar users as a guide.

\section*{Performance Metrics}

\begin{center}
\begin{tabular}{lp{9cm}}
\toprule
Metric & Definition \\
\midrule
Speedup & $S(p) = T_{\text{serial}} / T_{\text{parallel}}(p)$ \\
Efficiency & $E(p) = S(p) / p$ — fraction of ideal speedup \\
MAE & $\frac{1}{|T|}\sum_{(u,i)\in T}|\text{pred}(u,i) - \text{actual}(u,i)|$ \\
Sim checksum & Sum of all $N^2$ similarity values — cross-version sanity check \\
\bottomrule
\end{tabular}
\end{center}

\section*{Output Tags (Parsed by Benchmark Script)}

\begin{verbatim}
[Data]   Users: 1000 | Items: 1000 | Sparsity: 70% | Test ratings: 29866
[Timing] Similarity matrix  : 1.3549 s
[Timing] Prediction phase   : 7.0277 s
[Timing] Total (sim+pred)   : 8.3826 s
[Check]  Sim-matrix checksum: 942.387323
[Eval]   MAE on test set    : 1.2574
[Eval]   RMSE on test set   : 1.4579
\end{verbatim}
(Actual serial-baseline output, 1000x1000, on the RTX 3050 reference machine.)

\section*{Amdahl's Law — The Fundamental Limit}

\[
S_{\max}(p) = \frac{1}{f + (1-f)/p}
\quad\xrightarrow{p\to\infty}\quad \frac{1}{f}
\]

$f$ = fraction of time that \emph{cannot} be parallelised (data generation,
evaluation).  Here the \emph{algorithmic} serial fraction is tiny
($f \approx 0.3\%$), but on this 8-core laptop the achievable speedup is capped
by memory bandwidth and reduced all-core clocks, giving an \emph{effective}
$f \approx 3.8\%$ ($S_{\max} \approx 26\times$).  In practice the best CPU result
is $\approx 6.3\times$ at 8 workers, while the GPU (CUDA) reaches
$\approx 38\times$.

\section*{How the Files Relate}

\begin{center}
\begin{tabular}{ll}
\toprule
Source file & Topic covered in this guide \\
\midrule
\texttt{serial/serial\_recommender.c} & Chapter 1 \\
\texttt{openmp/openmp\_recommender.c} & Chapter 2 \\
\texttt{pthreads/pthreads\_recommender.c} & Chapter 3 \\
\texttt{mpi/mpi\_recommender.c} & Chapter 4 \\
\texttt{cuda/cuda\_recommender.cu} & Chapter 5 \\
\texttt{hybrid/hybrid\_recommender.c} & Chapter 6 \\
\bottomrule
\end{tabular}
\end{center}

% =============================================================================
\chapter{Serial Baseline \& Algorithm Foundation}
% =============================================================================

\section{Why a Serial Baseline?}

Before parallelising anything, every HPC project must have a correct serial
implementation serving three roles:

\begin{enumerate}
    \item \textbf{Correctness reference}: every parallel MAE must match within $\pm 0.001$
    \item \textbf{Speedup denominator}: $S = T_{\text{serial}} / T_{\text{parallel}}$
    \item \textbf{Algorithm clarity}: understand the logic before adding concurrency
\end{enumerate}

\section{The 5-Phase Algorithm}

\begin{center}
\begin{tabular}{lll}
\toprule
Phase & Operation & Complexity \\
\midrule
1 & Data Generation & $O(N \times M)$ \\
2 & User Means & $O(N \times M)$ \\
3 & Similarity Matrix & $O(N^2 \times M)$ \textbf{— Bottleneck} \\
4 & Predictions & $O(N^2 \times M)$ \\
5 & MAE Evaluation & $O(|T|)$ \\
\bottomrule
\end{tabular}
\end{center}

\section{Pearson Correlation Formula}

\[
\text{sim}(u,v) = \frac
{\sum_{i \in I_{uv}} (R_{u,i} - \bar{R}_u)(R_{v,i} - \bar{R}_v)}
{\sqrt{\sum_{i\in I_{uv}}(R_{u,i}-\bar{R}_u)^2}
 \times\sqrt{\sum_{i\in I_{uv}}(R_{v,i}-\bar{R}_v)^2}}
\]

where $I_{uv}$ = items rated by \emph{both} $u$ and $v$, and
$\bar{R}_u$ = mean rating of user $u$.

\textbf{Result range}: $[-1, +1]$. $+1$ = always agree, $0$ = unrelated,
$-1$ = always disagree.

\section{Prediction Formula}

\[
\text{pred}(u,i) = \bar{R}_u +
\frac{\sum_{k\in K}\text{sim}(u,k)\cdot(R_{k,i}-\bar{R}_k)}{\sum_{k\in K}\text{sim}(u,k)}
\]

$K$ = top-20 most similar neighbours who rated item $i$ with positive similarity.

\section{Key Constants}

\begin{lstlisting}
#define DEFAULT_USERS  1000    /* rows in ratings matrix */
#define DEFAULT_ITEMS  1000    /* columns */
#define SPARSITY       0.70f   /* 70% of entries are unrated */
#define TOP_K            20    /* max neighbours for prediction */
#define SEED             42    /* fixed RNG seed (same data every run) */
#define TEST_RATIO      0.10f  /* 10% held out for MAE */
\end{lstlisting}

% =============================================================================
\chapter{OpenMP — Shared-Memory Parallelism}
% =============================================================================

\section{Fork-Join Model}

OpenMP uses compiler directives (\texttt{\#pragma omp}) to automatically
manage a team of threads that share the same address space.

\section{Directives Used in This Project}

\begin{center}
\begin{tabular}{lll}
\toprule
Phase & Directive & Schedule \\
\midrule
2 (means) & \texttt{\#pragma omp parallel for} & \texttt{static} \\
3 (sim) & \texttt{\#pragma omp parallel for} & \texttt{dynamic, 4} \\
4 (pred) & \texttt{\#pragma omp parallel} + \texttt{for} & \texttt{dynamic, 2} \\
5 (MAE) & \texttt{\#pragma omp parallel for reduction(+:err)} & \texttt{static} \\
\bottomrule
\end{tabular}
\end{center}

\section{Schedule Comparison}

\begin{center}
\begin{tabular}{lp{6cm}l}
\toprule
Schedule & When to use & Overhead \\
\midrule
\texttt{static} & Equal work per iteration & None \\
\texttt{dynamic,C} & Unequal/unpredictable work & Queue lookup \\
\bottomrule
\end{tabular}
\end{center}

Phase 2 uses \texttt{static} (all rows scan exactly \texttt{N\_ITEMS} items).
Phases 3 \& 4 use \texttt{dynamic} (work varies with sparsity and triangle shape).

\section{Thread-Private Buffer Pattern (Phase 4)}

\begin{lstlisting}
#pragma omp parallel        /* fork threads */
{
    SimPair *nbrs = malloc(N_USERS * sizeof(SimPair));  /* private per thread */
    #pragma omp for schedule(dynamic, 2)
    for (int u = 0; u < N_USERS; u++) { /* distributed */ }
    free(nbrs);
}
\end{lstlisting}

\section{Why \texttt{omp\_get\_wtime()} not \texttt{clock()}?}

\texttt{clock()} sums CPU time across all threads ($4 \times 1\,\text{s} = 4\,\text{s}$
for 4 threads each working 1 second).  \texttt{omp\_get\_wtime()} measures
actual wall-clock elapsed time ($1\,\text{s}$).  Speedup requires wall-clock time.

% =============================================================================
\chapter{Pthreads — Manual Thread Management}
% =============================================================================

\section{Key Difference from OpenMP}

Pthreads requires explicit thread creation, work partitioning, and joining.
OpenMP hides all of this behind compiler directives.

\section{Core Pattern}

\begin{lstlisting}
/* Thread function */
void *thread_fn(void *arg) {
    ThreadArgs *a = (ThreadArgs *)arg;
    int start, end;
    work_range(N_USERS, a->tid, a->nthreads, &start, &end);
    /* ... process rows [start, end) ... */
    return NULL;
}

/* Launch in main() */
for (int t = 0; t < N_THREADS; t++) {
    args[t] = (ThreadArgs){t, N_THREADS};
    pthread_create(&threads[t], NULL, thread_fn, &args[t]);
}
for (int t = 0; t < N_THREADS; t++)
    pthread_join(threads[t], NULL);  /* wait = phase barrier */
\end{lstlisting}

\section{Ceiling Division for Work Partitioning}

\begin{lstlisting}
void work_range(int total, int tid, int nthreads, int *start, int *end) {
    int chunk = (total + nthreads - 1) / nthreads;
    *start = tid * chunk;
    *end   = min(*start + chunk, total);
}
\end{lstlisting}

Ensures all $N$ items are covered with no gaps and no overlaps even when
$N$ is not divisible by the thread count.

\section{No Mutex Needed}

Our code avoids mutexes entirely because each thread writes to a disjoint
range of indices.  Mutexes would serialise access and eliminate the speedup
benefit.

% =============================================================================
\chapter{MPI — Distributed-Memory Parallelism}
% =============================================================================

\section{Core Concept: Private Memory Spaces}

Each MPI rank is an independent process with its own private RAM.  Data
must be explicitly sent and received using MPI functions — no shared pointers.

\section{Initialisation}

\begin{lstlisting}
MPI_Init(&argc, &argv);
MPI_Comm_rank(MPI_COMM_WORLD, &mpi_rank);  /* 0-based rank ID */
MPI_Comm_size(MPI_COMM_WORLD, &mpi_size);  /* total ranks */
/* ... */
MPI_Finalize();   /* must be last MPI call */
\end{lstlisting}

\section{Key Collectives}

\begin{center}
\begin{tabular}{lp{9cm}}
\toprule
Call & Effect \\
\midrule
\texttt{MPI\_Allgatherv} &
  Each rank contributes a slice; every rank receives the complete assembled
  result.  Used for \texttt{user\_mean} and \texttt{sim\_matrix}. \\
\texttt{MPI\_Reduce} &
  Combines one value per rank (e.g.\ partial error sum) to rank 0.
  Used for MAE. \\
\texttt{MPI\_Barrier} &
  All ranks wait until every rank reaches the barrier.
  Used to bracket timed sections. \\
\bottomrule
\end{tabular}
\end{center}

\section{Parallelisation Per Phase}

\begin{center}
\begin{tabular}{lll}
\toprule
Phase & Local work & Communication \\
\midrule
1 (data gen) & Each rank runs \texttt{srand(42)} independently & None \\
2 (means) & Each rank computes its rows & \texttt{MPI\_Allgatherv} \\
3 (sim) & Each rank fills all columns for its rows & \texttt{MPI\_Allgatherv} \\
4 (pred) & Each rank predicts for its users only & None \\
5 (MAE) & Each rank sums its partial error & \texttt{MPI\_Reduce} \\
\bottomrule
\end{tabular}
\end{center}

% =============================================================================
\chapter{CUDA — GPU Massively Parallel Computing}
% =============================================================================

\section{CPU vs.\ GPU}

GPU has thousands of simple cores; ideal for data-parallel work where the same
operation applies to many independent elements.  Our $N^2$ independent Pearson
computations are a textbook GPU workload.

\section{Thread Hierarchy}

\textbf{Thread}: smallest unit, own registers, unique \texttt{threadIdx}\\
\textbf{Block}: up to 1024 threads sharing fast shared memory (\texttt{\_\_syncthreads})\\
\textbf{Grid}: all blocks for one kernel launch, scheduled across all SMs

\section{The Three Kernels}

\begin{center}
\begin{tabular}{lll}
\toprule
Kernel & Grid & Thread maps to \\
\midrule
\texttt{kernel\_user\_means} & 1D & One thread per user \\
\texttt{kernel\_similarity} & 2D & One thread per $(u,v)$ pair \\
\texttt{kernel\_predictions} & 2D & One thread per $(u,\text{item})$ pair \\
\bottomrule
\end{tabular}
\end{center}

\section{Memory Flow}

\begin{center}
CPU RAM $\xrightarrow{\texttt{cudaMemcpy (H2D)}}$ GPU DRAM
$\xrightarrow{\text{kernel}}$ GPU DRAM
$\xrightarrow{\texttt{cudaMemcpy (D2H)}}$ CPU RAM
\end{center}

\section{Timing GPU Code}

Use CUDA events (not \texttt{clock\_gettime}) because kernel launches are
\emph{asynchronous}:
\begin{lstlisting}
cudaEventRecord(ev_start);
kernel<<<grid,block>>>(...);
cudaEventRecord(ev_stop);
cudaEventSynchronize(ev_stop);
float ms; cudaEventElapsedTime(&ms, ev_start, ev_stop);
\end{lstlisting}

% =============================================================================
\chapter{Hybrid MPI+OpenMP — Two-Level Parallelism}
% =============================================================================

\section{Motivation}

Real clusters have multi-core nodes connected by a network.
\begin{itemize}
    \item \textbf{Pure MPI}: intra-node communication is wasteful (cores share RAM)
    \item \textbf{Pure OpenMP}: cannot cross node boundaries
    \item \textbf{Hybrid}: MPI between nodes, OpenMP within nodes
\end{itemize}

\section{Thread-Safe MPI Initialisation}

\begin{lstlisting}
int provided;
MPI_Init_thread(&argc, &argv, MPI_THREAD_FUNNELED, &provided);
/* MPI_THREAD_FUNNELED: MPI only called from main thread */
\end{lstlisting}

\section{Two-Level Work Division}

\begin{center}
\begin{tabular}{lll}
\toprule
Level & Technology & Grain \\
\midrule
Coarse & MPI & Rows assigned to different ranks \\
Fine & OpenMP & Threads share rows within a rank \\
\bottomrule
\end{tabular}
\end{center}

Total workers = $P \times T$, MPI messages = $P$ (not $P \times T$).

\section{Two-Level Reduction (Phase 5)}

\begin{enumerate}
    \item OpenMP \texttt{reduction(+:local\_err, local\_cnt)}: private per-thread copies merged at barrier
    \item \texttt{MPI\_Reduce}: per-rank totals summed to rank 0
\end{enumerate}

\section{Benchmarked Configurations}

\begin{center}
\begin{tabular}{lll}
\toprule
Config & $P \times T$ & Notes \\
\midrule
Hybrid\_2x4 & 2 ranks, 4 threads & Balanced \\
Hybrid\_4x2 & 4 ranks, 2 threads & More MPI \\
Hybrid\_8x1 & 8 ranks, 1 thread & Equivalent to pure MPI \\
Hybrid\_1x8 & 1 rank, 8 threads & Equivalent to pure OpenMP \\
\bottomrule
\end{tabular}
\end{center}

% =============================================================================
\chapter*{Summary Comparison}
\addcontentsline{toc}{chapter}{Summary Comparison}
% =============================================================================

\begin{center}
\begin{tabular}{lllll}
\toprule
Version & Memory & Parallelism & Best for & Communication \\
\midrule
Serial & Shared & None & Baseline & None \\
OpenMP & Shared & Fork-join & Single machine & None \\
Pthreads & Shared & Manual threads & Full control & None \\
MPI & Distributed & Processes & Multi-node & Messages \\
CUDA & Host+GPU & 1000s of threads & Data-parallel & PCIe transfer \\
Hybrid & Distributed & MPI$\times$OMP & Large clusters & MPI only \\
\bottomrule
\end{tabular}
\end{center}

\vspace{1em}
\begin{center}
\textbf{Measured Speedup Order} (this project; 8 CPU workers, CUDA on RTX 3050)\\[0.5em]
\texttt{CUDA} ($\times38$) $>$ \texttt{OpenMP} ($\times6.3$) $>$ \texttt{Pthreads} ($\times5.7$)
$>$ \texttt{MPI} ($\times5.1$) $>$ \texttt{Hybrid\,8$\times$1} ($\times4.9$) $>$ \texttt{Serial}
\end{center}

\end{document}
